\section{RQ2: Deep Analysis for Prefix Hint}

\subsection{RQ2.1 Hint-conditioned training signal 是否有效}

这个问题先不看 family 内部细节，只问一件事：当训练过程真正看到 prefix hint 之后，信号是否会转化成可观测的性能变化，而不只是把 \texttt{rule-only} 换个包装。

\begin{table}[H]
\centering
\small
\setlength{\tabcolsep}{4pt}
\begin{tabular}{p{3.9cm} p{2.3cm} c c c c}
\toprule
Variant & Best ckpt & NDCG@10 & HR@10 & NDCG@50 & HR@50 \\
\midrule
SFT495 & checkpoint-495 & 0.0823 & 0.1094 & 0.0985 & 0.1844 \\
rule-only rerun & checkpoint-2997 & 0.0960 & 0.1179 & 0.1070 & 0.1681 \\
dynamic gather-fix & checkpoint-2997 & 0.0936 & 0.1169 & 0.1083 & 0.1855 \\
fixed taskfix & checkpoint-2997 & 0.0931 & 0.1189 & 0.1094 & 0.1941 \\
fixed taskfix sid-only & checkpoint-2652 & 0.0945 & 0.1205 & 0.1103 & 0.1935 \\
\bottomrule
\end{tabular}
\caption{hint-conditioned training signal 的第一层证据。}
\end{table}

这张表背后的判断很直接：

\begin{itemize}
\item \texttt{rule-only} 把 \texttt{NDCG@10} 推高了，但 \texttt{HR@50} 从 \texttt{0.1844} 掉到 \texttt{0.1681}，说明无 hint 的 exact reward 仍然主要在做 top-10 极值化。
\item \texttt{dynamic gather-fix} 在只比 \texttt{rule-only} 低 \texttt{0.0024} 个 \texttt{NDCG@10} 的同时，把 \texttt{HR@50} 拉回到 \texttt{0.1855}，已经超过 SFT；这就是“hint-conditioned signal 有效”的最直接证据。
\item 两条 fixed 线继续把右上角往外推，尤其是 \texttt{fixed taskfix sid-only}，说明 prefix hint 的收益并不只发生在 adaptive 侧。
\end{itemize}

\subsection{RQ2.2 Prefix Hint Strategy 的 2x2 对比}

这里把 prefix-hint strategy 压缩成一个明确的 \texttt{2x2}：

\begin{table}[H]
\centering
\small
\begin{tabular}{lcc}
\toprule
 & task-aware prefix (ours) & first-token prefix \\
\midrule
fixed hint & \texttt{fixed taskfix} & \texttt{fixed taskfix sid-only} \\
adaptive hint & \texttt{dynamic gather-fix} & \texttt{dynamic sid-only} \\
\bottomrule
\end{tabular}
\caption{RQ2 中的 prefix-hint 2x2。}
\end{table}

\begin{table}[H]
\centering
\small
\setlength{\tabcolsep}{4pt}
\begin{tabular}{p{3.9cm} p{2.3cm} c c c c}
\toprule
Variant & Best ckpt & NDCG@10 & HR@10 & NDCG@50 & HR@50 \\
\midrule
fixed taskfix (ours) & checkpoint-2997 & 0.0931 & 0.1189 & 0.1094 & 0.1941 \\
dynamic gather-fix (adaptive hinting) & checkpoint-2997 & 0.0936 & 0.1169 & 0.1083 & 0.1855 \\
fixed taskfix sid-only (fixed first-token) & checkpoint-2652 & 0.0945 & 0.1205 & 0.1103 & 0.1935 \\
dynamic sid-only (adaptive first-token) & checkpoint-2394 & 0.0921 & 0.1155 & 0.1065 & 0.1830 \\
\bottomrule
\end{tabular}
\caption{prefix-hint 2x2 在 best checkpoint 上的对照。}
\end{table}

\begin{figure}[H]
\centering
\includegraphics[width=0.96\textwidth]{prefix_hint_strategy_2x2_curves.png}
\caption{prefix-hint 2x2 的完整 epoch 曲线。}
\end{figure}

这组对照可以回答两个更细的问题：

\begin{itemize}
\item 在 adaptive 行里，\texttt{dynamic gather-fix} 四个主指标同时强于 \texttt{dynamic sid-only}，说明 adaptive hinting 本身仍然需要 task-aware prefix，而不是退化成 first-token 就够。
\item 在 fixed 行里，\texttt{fixed taskfix sid-only} 把 \texttt{NDCG@10} 从 \texttt{0.0931} 提到 \texttt{0.0945}，而 \texttt{fixed taskfix} 保住了略高的 \texttt{HR@50}；因此 fixed 侧更像是在 top-10 与 coverage 之间做形状调整，而不是单边支配。
\item 跨行看，当前 best-point 意义下 fixed 侧整体仍强于 adaptive 侧，因此本文后面的 loss/CE 讨论都优先放在 fixed family 上展开。
\end{itemize}

\subsection{RQ2.3 Prefix Budget 敏感性：max1}

如果把 adaptive hint 的在线 budget 从 \texttt{max depth=3} 收紧到 \texttt{1}，可以更直接地回答“hint signal 是否真的来自更深的 prefix”。

\begin{table}[H]
\centering
\small
\setlength{\tabcolsep}{4pt}
\begin{tabular}{p{3.8cm} p{2.4cm} c c c}
\toprule
Variant & Best ckpt & Epoch & NDCG@10 & HR@50 \\
\midrule
rule-only rerun & checkpoint-2997 & 1.802 & 0.0960 & 0.1681 \\
dynamic gather-fix & checkpoint-2997 & 1.802 & 0.0936 & 0.1855 \\
dynamic max1 & checkpoint-1332 & 0.801 & 0.0934 & 0.1905 \\
fixed taskfix sid-only & checkpoint-2652 & 2.000 & 0.0945 & 0.1935 \\
\bottomrule
\end{tabular}
\caption{max1 说明“只给一个 prefix token”依然有信号，但它更像 early-stop candidate。}
\end{table}

\begin{figure}[H]
\centering
\includegraphics[width=0.96\textwidth]{max1-ablation-epoch-curves.png}
\caption{max1 ablation 曲线。}
\end{figure}

\begin{figure}[H]
\centering
\includegraphics[width=0.96\textwidth]{max1-vs-fixed-epoch-curves.png}
\caption{max1 与 dynamic / fixed 参考线的 focused comparison。}
\end{figure}

\begin{itemize}
\item \texttt{max1} 没有退化成 plain \texttt{rule-only}，因为它保住了显著更高的 coverage。
\item 但完整轨迹表明，它的优势主要停在前中段；和 \texttt{dynamic gather-fix} 相比，它更像一个 budget-constrained early-stop 变体，而不是新的 long-run default。
\item 这也反过来说明 RQ2 的主结论不是“任意 prefix 都行”，而是：prefix hint 确实有效，但深度、task awareness 和 fixed/adaptive 交互方式都会改变最终 trade-off 的形状。
\end{itemize}
